\chapter*{Preface}

\epigraph{Mathematics is the art of giving the same name to different things.}{\textit{Henri Poincaré}}

\noindent This book is something of a Trojan horse.

On its surface, it presents itself as a complement to existing advanced undergraduate or introductory graduate texts in quantitative macroeconomics and finance---works like \textit{Dynamic Economics} by Adda and Cooper or \textit{Recursive Macroeconomic Theory} by Ljungqvist and Sargent. In some respects it is less technical than these references, and certainly less comprehensive. But it pursues a different aim: to expose the common structure that underlies an enormous range of applied problems in macroeconomics, finance, and international trade. The organizing principle is the constant elasticity of substitution---the CES operator and its many manifestations---which appears again and again as the workhorse aggregator in models of heterogeneous agents, firms, and intermediaries.

The intended audience is threefold. First, advanced undergraduates and beginning graduate students who want an accessible entry point into dynamic macroeconomics and finance. Second, researchers in macro, finance, and trade who may benefit from seeing their problems mapped into a common structure with potential unexpected connections. Third---and perhaps most honestly---myself. As a researcher, I have long felt that many problems in our field are reformulations of the same underlying forces, yet we often derive identical formulas from scratch, paper after paper. I wanted a unified treatment that would spare me (and others if they wish) of this repetition, while also serving as a glossary of mechanisms for teaching and for my own reference. Maybe it is just a training ground to be faster and more silent in a seminar. 

The scope is deliberately narrow in certain dimensions. I have left out financial constraints, information frictions, the nominal-real dichotomy, and most forms of non-convexity---no S-s bands, no lumpy adjustment. No heterogeneity in preferences. This is not because these features are unimportant; quite the opposite. But the essence of the problems I wish to illuminate does not change fundamentally when we add these layers, and their inclusion would obscure the structural isomorphisms I want to emphasize. The book focuses instead on what is analytically tractable and computationally portable: settings where we can derive closed-form aggregators and see clearly how idiosyncratic shocks, preference parameters, and equilibrium forces interact.

Several developments in the profession inspired this project. First, I have been struck by how much recent work explores the interaction between risk, flexibility, and intertemporal substitution---themes that cut across asset pricing, business cycles, trade, and firm dynamics. Second, I have observed a growing reliance on computational brute force, often at the expense of understanding what forces in a model actually drive the results. Papers frequently leap to numerical solutions without first asking what a representative-agent version of their economy would predict, or which parameters matter most for the mechanism at hand. We forget about substitution and wealth effects, key decompositions to understand policy implications and modeling assumptions. One purpose of this book is to provide precisely that preliminary map: a glossary of forces that helps researchers (and students) understand what their models are doing before they write a single line of code.

The intellectual debts are vast. The concepts themselves are embarrasingly old---CES aggregation, Bellman equations, distributional dynamics---and I claim no originality. What I hope is new is the synthesis: the explicit mapping of problems across subfields, the repeated use of the same formulas in different guises, and the emphasis on building intuition to complement computation. I have been particularly influenced by the lecture notes of Thomas Sargent and John Stachurski, whose commitment to making quantitative methods accessible has shaped how I think about pedagogy. But where their work emphasizes breadth and computational implementation, mine emphasizes depth and analytical structure---a bridge, I hope, rather than a substitute.

The Trojan horse brought warriors hidden inside a gift. This book brings, hidden inside a standard advanced undegraduate material, a perspective. The warriors, if you will, are the recurring structures. However, the warriors at our service.

\vspace{1cm}
\hfill\textit{S.B.}

\hfill\textit{Los Angeles, 2026}